\subsection{GPU 基础：峰值算力之外，先看数据能否喂到计算单元}

slides 58--63 建立系统部分的最低知识框架。GPU 用大量线程执行相同指令并擅长矩阵乘法，因此适合高吞吐训练；但计算单元增长速度远快于内存和通信，真正瓶颈常是把权重、activation 与梯度搬到合适位置。memory hierarchy 越靠近计算核心越快、容量越小，越远越慢、容量越大。Model FLOP Utilization（MFU）把观测吞吐与硬件理论峰值相除，用于判断训练是否有效利用矩阵算力。

\lecturefigure{slides-images/slide-058.png}{官方 slide 58：GPU 稀缺、物理通信限制与资源分配}
\lecturefigure{slides-images/slide-059.png}{官方 slide 59：GPU 的线程级大规模并行}
\lecturefigure{slides-images/slide-060.png}{官方 slide 60：专用矩阵核心带来的高吞吐}
\lecturefigure{slides-images/slide-061.png}{官方 slide 61：计算增长快于 memory/communication}
\lecturefigure{slides-images/slide-062.png}{官方 slide 62：靠近核心更快但更小的 memory hierarchy}
\lecturefigure{slides-images/slide-063.png}{官方 slide 63：MFU 定义与 50\% 的工程直觉}

\begin{knowledgebox}{术语消化：DRAM、HBM、SRAM 与 MFU}
\textbf{DRAM（Dynamic Random-Access Memory）}是动态随机存取存储器；GPU 显存通常采用高带宽封装的 \textbf{HBM（High Bandwidth Memory）}，容量大但离计算核心较远。\textbf{SRAM（Static Random-Access Memory）}用于片上 cache/shared memory，速度快但容量小。数据从 HBM 反复读写会浪费带宽和能量，因此 kernel 会尽量在 SRAM/寄存器中复用。MFU 定义为实际有效模型 FLOPs 吞吐除以理论峰值；它不是 GPU utilization 百分比，也不包含所有非模型操作的价值判断。
\end{knowledgebox}

\begin{knowledgebox}{读图：DataMovement 图在说明什么}
slide 61 将 BERT 算子中的 MatMul、Activation 与 DataMovement 分开，目的是说明 wall-clock 时间不只由 FLOPs 决定。一个算子 FLOPs 很少，却可能因为每次都读写大张量而耗时；另一个矩阵乘法 FLOPs 很高，但 Tensor Core 能高效完成。读这类图时先区分“计算量占比”和“执行时间占比”，再问数据是否被复用、是否触发跨卡通信。不能从 FLOPs 低直接推出优化价值低。
\end{knowledgebox}

\begin{warningbox}{50\% MFU 为什么已经很好}
\textbf{老师强调：}约 01:36，讲者说 50\% MFU 已处于非常好的状态，因为真实训练包含通信、非矩阵算子、数据加载、同步和数值稳定处理。理论峰值通常假设理想矩阵形状与持续满载。MFU 低需要定位瓶颈，但 MFU 高也不代表端到端系统最优：若为了堆大 batch 提高 MFU，却降低样本效率或增加延迟，最终训练成本仍可能更差。
\end{warningbox}

Agent RL 会进一步放大系统不规则性：不同轨迹长度不同，工具响应时间不同，部分环境需要网络或沙箱，GPU 很难像预训练那样持续处理整齐的大 batch。因此课程在 RL 部分提到暂停长尾 rollout 和并发环境服务，与这里的 memory/compute 分析是一条主线：优化目标不是让每个组件单独满载，而是提高单位时间获得的有效学习信号。

